\documentclass{article}
\usepackage{amsmath,amssymb,mathtools}
\begin{document}

\section*{Higham Chapter 6 Problem 6.5: unitarily invariant norm blocker}

I am formalizing Higham, \emph{Accuracy and Stability of Numerical
Algorithms}, 2nd ed., Chapter 6, Problem 6.5.  The local Lean repository already
has a checked complex Frobenius/operator-2 proof of
\[
  \|ABC\|_F \le \|A\|_2\,\|B\|_F\,\|C\|_2 .
\]
The remaining source sentence is the arbitrary unitarily invariant norm
extension:
\[
  |||ABC||| \le \|A\|_2\,|||B|||\,\|C\|_2
\]
for every unitarily invariant matrix norm.

\subsection*{Local formalization constraints}

The local chapter file has:
\begin{itemize}
\item finite complex matrices as functions \texttt{Fin m -> Fin n -> Complex};
\item matrix multiplication, adjoint, Frobenius norm, operator 2-norm;
\item rectangular SVD/factorization and local singular values;
\item no general majorization/Fan dominance API in Mathlib or the project.
\end{itemize}

Another project file already has a real rectangular
\texttt{UnitaryInvariantRectNormLike} abstraction with orthogonal left/right
invariance only.  It does \emph{not} prove the ideal inequalities.  A safe local
API could be an ``operator ideal norm'' family with fields
\[
  |||LX||| \le \|L\|_2\,|||X|||,\qquad
  |||XR||| \le |||X|||\,\|R\|_2,
\]
from which the triple product theorem is immediate.  But this does not prove
that every unitarily invariant norm has those fields.

\subsection*{Questions}

\begin{enumerate}
\item What is the shortest rigorous finite-dimensional proof that every
unitarily invariant matrix norm satisfies
\[
  |||LX||| \le \|L\|_2\,|||X|||
\]
and the right-hand analogue?  Please give proof-assistant-sized lemmas, not just
a citation.
\item Is the standard route necessarily Fan dominance / weak majorization of
singular values?  If so, list the exact intermediate theorems needed:
singular-value inequality, monotonicity of unitarily invariant norms under weak
majorization, symmetric gauge representation, etc.
\item Is there a shorter proof using contraction dilation or polar/SVD
factorization that avoids building a full majorization theory, especially in
finite dimensions and for rectangular matrices?  If yes, spell out the required
dimension-padding and norm-extension assumptions.
\item For a Lean core pass, is it mathematically honest to close the source row
with a theorem over an explicit ``unitarily invariant operator ideal norm''
interface and record Fan dominance as the missing foundation for deriving the
interface from bare unitary invariance?  Or would that still be hiding the main
problem?
\item Please separate theorem facts from suggested formalization shortcuts.
\end{enumerate}

\end{document}
